% =====================================================================
%  c00b.tex —— 致谢 / 摘要 / 译者说明（Lead 提供）
%  原文：印刷 p.iv（致谢）、印刷 p.v（摘要）
% =====================================================================

\chapter*{致谢}
\addcontentsline{toc}{chapter}{致谢}
\markboth{致谢}{致谢}

% [原文 p.iv]

我要感谢我的导师 Anton Kapustin：本论文的核心思想来自他；我从他那里学到了极多的东西；
在写作过程中，他以指导与支持使我始终保持在正轨上。同样感谢 Sergei Gukov 的指导与支持。

非常感谢 Carol Silberstein 和理论组的各位成员，尤其是 Ketan Vyas、Paul Cook、Jie Yang、
Sakura Schafer-Namecki、Tudor Dimofte、Miguel Bandres、Itamar Yaakov、Brian Willett 与
Denis Bashkirov，感谢他们的讨论与支持。感谢 Lee Coleman 与 Cristina Dezan 对生活以及更多
事情的大量洞见。我感谢 Jack Kent Cooke 基金会的慷慨资助。

诚挚感谢我的健身伙伴 Chris Wegg 与 Nate Bode，以及我的朋友 Lucia Cordeiro-Drever、Esi Alizadeh、
Laura Book、Tristan Smith 与 Hernan Garcia，感谢他们在我在 Caltech 期间给予的支持与启发。
深深的爱与感激献给我的父母，感谢他们不知疲倦的支持与理解。

最后，向 Camilla Berretta 致以我最深的谢意——
\begin{quote}
\itshape
So are you to my thoughts as food to life\\
Or as sweet-season'd showers are to the ground.
\end{quote}
（你之于我的思念，犹如食物之于生命，又如甘霖之于大地。）
% [Tutor Note: 这两行出自莎士比亚十四行诗第 75 首（Sonnet 75），原文为 "So are you to my thoughts as food to life, / Or as sweet-season'd showers are to the ground"；作者引用以表谢意，此处保留英文并附中译。]

\chapter*{摘要}
\addcontentsline{toc}{chapter}{摘要}
\markboth{摘要}{摘要}

% [原文 p.v]

在 A. Kapustin 与 E. Witten 的开创性工作 [1] 中，数论中的几何朗兰兹纲领被证明与紧化在
一个黎曼面上的 GL-扭转 $\cN=4$ 超杨-米尔斯理论的对偶性密切相关。本论文推广
Kapustin--Witten 的工作：我们研究 GL-扭转理论在圆上紧化到三维的情形（对扭转参数 $t$ 的
若干取值）。通过考察三维描述中的边界条件，我们对 GL-扭转理论的余维二面算子进行分类，
推广了 S. Gukov 与 E. Witten [2] 所研究的那类面算子。对于 $t=\ii$，我们用模范畴给出
面算子 2-范畴的完整描述，此外还确定了包含 Wilson 线作为特殊对象的线算子幺半范畴。
对于 $t=1$ 与 $t=0$，我们讨论阿贝尔情形下的面算子与线算子。

我们还通过分析 $\cN=4$ 的另一种扭转版本——Vafa-Witten 理论——来推广 Kapustin--Witten 的工作。
在引入一个新的四维拓扑规范理论——带规范场的 4d A-模型——之后，我们把 Vafa-Witten 理论
定位为其特殊情形。我们讨论了 Vafa-Witten 理论在圆上与在黎曼面上的紧化。
本论文及附录中还研究了若干新的二维与三维拓扑规范理论。

在与本论文主线无关的工作中，我们在结尾对具有不同超对称量的二维 $\sigma$-模型中的
余维一拓扑缺陷进行了分类。

% [Tutor Note: 摘要中“紧化在一个黎曼面上”（compactified on a Riemann surface）与后文“在圆上紧化”
%  对举：前者保留两个维度（$\Sigma\times C$ 中的 $C$ 收缩），后者把时空的一个方向紧化到 $S^1$，
%  从而得到 3d 有效理论；这正是本论文相对 Kapustin--Witten 的主要推广方向。]

\chapter*{译者说明}
\addcontentsline{toc}{chapter}{译者说明}
\markboth{译者说明}{译者说明}

\para{底本。}
Kevin Setter, \emph{Topological Quantum Field Theory and the Geometric Langlands Correspondence},
Thesis by, California Institute of Technology, Pasadena, California, 2013（2012 年 8 月 24 日答辩；
导师 Anton Kapustin）。原文 PDF 共 169 页；本译本正文覆盖原文印刷 pp.1--154（9 章 + 附录 A--D），
参考文献（印刷 pp.155--159，64 条）保留英文原貌。

\para{页码映射。}
原文 PDF 页号 $=$ 印刷页码 $+ 10$（PDF p.11 的页脚为“1”，PDF p.169 的页脚为“159”）。
本译本每页内容前用注释标注原文页码：\verb|% [原文 p.N]|，便于与原文逐页对照。

\para{结构。}
第 1 章引言与概述；第 2 章四维规范理论的对偶；第 3 章 Vafa-Witten 理论作为带规范场的 4d A-模型；
第 4 章拓扑场论、范畴与 2-范畴；第 5 章在圆上紧化；第 6 章 4d TFT 的面算子；第 7 章修正 A-模型；
第 8 章 Vafa-Witten 理论在亏格 $g\ge 2$ 黎曼面上的紧化；第 9 章二维 $\sigma$-模型的缺陷；
附录 A（2d TFT）、附录 B（3d TFT）、附录 C（相伴纤维丛与等变上同调）、附录 D（第 9 章命题的证明）。

\para{体例。}
\begin{itemize}[leftmargin=1.6em,itemsep=1pt,topsep=2pt]
\item 编号公式严格照抄原文编号（\verb|\tag{N.M}|），与原文一一对应，不增、不删、不重排；
      原文中没有编号的公式一律不编号。正文引用公式时写作“式 (2.3)”这样的纯文本。
\item 原文 10 张图（图 1.1、2.1、4.1--4.7、6.1）均按原文页面图像 400 dpi 截取，图注由译者译出。
\item 原文的方框公式（如式 (1.1)--(1.3)）按原样用带边框的排版复现。
\item 定理、定义、命题、习题等环境的编号一律手工照抄原文（如“定义 2”“命题 7”），
      不按 LaTeX 计数器自动编号。
\item 术语以附录后的术语对照表为准；人名、机构名、模型名（Kapustin、Witten、Gukov、Hitchin、
      Vafa-Witten、Nahm、Rozansky--Witten 等）保留原文。
\item 译者补充的解释、省略论证的补全、以及发现的原文笔误，全部写在
      \verb|% [Tutor Note: ...]| 注释中，不进入排版；其索引见随书交付的 \texttt{译注汇总.md}。
\end{itemize}

\para{编译。}
本译本用 XeLaTeX 编译三遍（源码与一键脚本见交付目录）：中文字体使用 Windows 自带的
SimSun/SimHei/SimKai 文件路径，不依赖 fontconfig。

\para{已知偏差。}
（1）原文页码、行宽与断行不可能逐一复刻，译本按 A4 版式重新排版，内容与公式编号与原文一一对应；
（2）参考文献只做编码层面的整理（如把 PDF 文字层中的连字符断行、连字 fi、变音符号还原），不改条目内容；
（3）原文个别笔误、可疑之处在批注中标注，正文照译。
